\documentclass[12pt,a4paper]{report}

% ============================================================
% PACKAGES
% ============================================================
\usepackage[top=25.4mm, left=37mm, bottom=32mm, right=25.4mm, bindingoffset=0mm]{geometry}
\usepackage{times}
\usepackage{setspace}
\usepackage{titlesec}
\usepackage{fancyhdr}
\usepackage{booktabs}
\usepackage{array}
\usepackage{longtable}
\usepackage{graphicx}
\usepackage{listings}
\usepackage{url}
\usepackage{hyperref}
\usepackage{indentfirst}
\usepackage{caption}
\usepackage{tocloft}
\usepackage{enumitem}
\usepackage{multirow}

% ============================================================
% FORMATTING SETTINGS
% ============================================================

% 1.5 line spacing
\onehalfspacing

% Paragraph indentation (5 spaces)
\setlength{\parindent}{2.5em}

% Footer: SITS, Narhe, Pune - 41 (B.E. Computer Engineering) - 10pt centered
\pagestyle{fancy}
\fancyhf{}
\fancyfoot[C]{\fontsize{10}{12}\selectfont SITS, Narhe, Pune - 41 (B.E. Computer Engineering)}
\fancyfoot[R]{\thepage}
\renewcommand{\headrulewidth}{0pt}
\renewcommand{\footrulewidth}{0pt}

% Apply footer to plain pages (chapter starts)
\fancypagestyle{plain}{
  \fancyhf{}
  \fancyfoot[C]{\fontsize{10}{12}\selectfont SITS, Narhe, Pune - 41 (B.E. Computer Engineering)}
  \fancyfoot[R]{\thepage}
  \renewcommand{\headrulewidth}{0pt}
}

% Chapter title: 14pt BOLD UPPERCASE CENTERED
\titleformat{\chapter}[display]
  {\normalfont\fontsize{14}{16}\bfseries\centering}
  {\MakeUppercase{\chaptertitlename}\ \thechapter}{0pt}
  {\fontsize{14}{16}\bfseries\MakeUppercase}
\titlespacing*{\chapter}{0pt}{-20pt}{20pt}

% Section: 12pt bold UPPERCASE left-aligned
\titleformat{\section}
  {\normalfont\fontsize{12}{14}\bfseries}
  {\thesection}{1em}{\MakeUppercase}
\titlespacing*{\section}{0pt}{24pt}{12pt}

% Subsection: 12pt bold Title Case left-aligned
\titleformat{\subsection}
  {\normalfont\fontsize{12}{14}\bfseries}
  {\thesubsection}{1em}{}
\titlespacing*{\subsection}{0pt}{12pt}{6pt}

% Code listings
\lstset{
  basicstyle=\ttfamily\small,
  breaklines=true,
  frame=single,
  language=SQL,
  keywordstyle=\bfseries,
}

% Table/Figure captions - 12pt
\captionsetup{font={normalsize}, labelfont={bf}}

% Hyperlinks
\hypersetup{
  colorlinks=true,
  linkcolor=black,
  citecolor=black,
  urlcolor=blue,
}

\begin{document}

% ============================================================
% COVER PAGE
% ============================================================
\begin{titlepage}
\begin{center}

\vspace*{2\baselineskip}

{\fontsize{12}{14}\selectfont \textbf{A PROJECT REPORT ON}}

\vspace{1.5\baselineskip}

{\fontsize{16}{18}\selectfont \textbf{RIVETDB: A CONCURRENT, FEATURE-RICH IN-MEMORY DATABASE IN RUST WITH BUILT-IN TIME-SERIES, SQL QUERY, AND MULTI-TENANCY SUPPORT}}

\vspace{1.5\baselineskip}

{\fontsize{12}{14}\selectfont \textbf{SUBMITTED TO THE SAVITRIBAI PHULE PUNE UNIVERSITY,\\
IN THE PARTIAL FULFILLMENT OF REQUIREMENTS\\
FOR THE AWARD OF THE DEGREE}}

\vspace{0.5\baselineskip}

{\fontsize{12}{14}\selectfont \textbf{OF}}

\vspace{1\baselineskip}

{\fontsize{16}{18}\selectfont \textbf{BACHELOR OF ENGINEERING (COMPUTER ENGINEERING)}}

\vspace{1.5\baselineskip}

{\fontsize{14}{16}\selectfont \textbf{BY}}

\vspace{1\baselineskip}

{\fontsize{14}{16}\selectfont \textbf{
\begin{tabular}{l}
AYAN SHAIKH \\
RAHUL POLAS \\
JATIN SHIMPI \\
GAURI SHINDE \\
\end{tabular}
}}

\vspace{1.5\baselineskip}

{\fontsize{16}{18}\selectfont \textbf{DEPARTMENT OF COMPUTER ENGINEERING}}

\vspace{0.5\baselineskip}

{\fontsize{14}{16}\selectfont \textbf{SINHGAD TECHNICAL EDUCATION SOCIETY'S\\
SINHGAD INSTITUTE OF TECHNOLOGY \& SCIENCE}}

\vspace{0.5\baselineskip}

{\fontsize{12}{14}\selectfont \textbf{NARHE (AMBEGAON), OFF WESTERLY BY PASS PUNE MUMBAI EXPRESSWAY}}

\vspace{0.5\baselineskip}

{\fontsize{14}{16}\selectfont \textbf{NARHE, PUNE - 411 041}}

\vspace{1.5\baselineskip}

{\fontsize{14}{16}\selectfont 2024--2025}

\end{center}
\end{titlepage}

% ============================================================
% CERTIFICATE
% ============================================================
\newpage
\thispagestyle{plain}
\begin{center}

\vspace*{2\baselineskip}

{\fontsize{16}{18}\selectfont \textbf{CERTIFICATE}}

\vspace{1.5\baselineskip}

\end{center}

\fontsize{12}{14}\selectfont
This is to certify that the Project report entitled

\vspace{0.5\baselineskip}

\begin{center}
{\fontsize{12}{14}\selectfont \textbf{``RIVETDB: A CONCURRENT, FEATURE-RICH IN-MEMORY DATABASE IN RUST WITH BUILT-IN TIME-SERIES, SQL QUERY, AND MULTI-TENANCY SUPPORT''}}
\end{center}

\vspace{1\baselineskip}

Submitted by

\vspace{0.5\baselineskip}

\begin{center}
\begin{tabular}{l}
Ayan Shaikh \\
Rahul Polas \\
Jatin Shimpi \\
Gauri Shinde \\
\end{tabular}
\end{center}

\vspace{0.5\baselineskip}

is a bonafide work carried out by them under the Guidance of Prof. Swati B. Gholap and it is approved for the partial fulfilment of the requirement of the Savitribai Phule Pune University for the award of the Bachelor's Degree of Engineering (Computer Engineering).

\vspace{0.5\baselineskip}

This project report has not been earlier submitted to any other institute or university for the award of any degree or diploma.

\vspace{3\baselineskip}

\begin{tabular}{p{4cm}p{5cm}p{4cm}}
\centering \rule{3.5cm}{0.4pt} & \centering \rule{4cm}{0.4pt} & \centering\arraybackslash \rule{3.5cm}{0.4pt} \\
\centering Guide & \centering Project Coordinator & \centering\arraybackslash HOD Comp. Engg. \\
\end{tabular}

\vspace{2\baselineskip}

\begin{center}
\rule{4cm}{0.4pt}\\
Principal, SITS
\end{center}

\vspace{1\baselineskip}

\textbf{Department of Computer Engineering}\\
\textbf{Sinhgad Institute of Technology \& Science}

\vspace{0.5\baselineskip}

Place: Pune\\
Date: \rule{3cm}{0.4pt}

% ============================================================
% ACKNOWLEDGEMENT
% ============================================================
\newpage
\thispagestyle{plain}
\begin{center}
\vspace*{1\baselineskip}
{\fontsize{14}{16}\selectfont \textbf{ACKNOWLEDGEMENT}}
\vspace{1.5\baselineskip}
\end{center}

We would like to express our sincere gratitude to our guide \textbf{Prof. Swati B. Gholap}, Department of Computer Engineering, Sinhgad Institute of Technology \& Science, Narhe, Pune, for her valuable guidance, constant encouragement, and cooperation throughout the development of this project.

We are also grateful to \textbf{Dr. [HOD Name]}, Head of the Department of Computer Engineering, for providing the necessary facilities and support.

We extend our heartfelt thanks to \textbf{Dr. [Principal Name]}, Principal, Sinhgad Institute of Technology \& Science, for providing excellent infrastructure and a conducive environment for project work.

We are thankful to all the teaching and non-teaching staff of the Department of Computer Engineering for their help and cooperation.

Finally, we would like to thank our families and friends for their constant support and motivation throughout this journey.

\vspace{4\baselineskip}

\begin{flushright}
Ayan Shaikh\\
Rahul Polas\\
Jatin Shimpi\\
Gauri Shinde\\
\end{flushright}

% ============================================================
% ABSTRACT
% ============================================================
\newpage
\thispagestyle{plain}
\begin{center}
{\fontsize{12}{14}\selectfont \textbf{Abstract}}
\end{center}
In-memory databases are essential for modern applications requiring sub-millisecond latency, yet existing systems present a fundamental trade-off: Redis offers feature richness through a single-threaded model that limits scalability, while Memcached provides multi-threaded throughput at the expense of functionality. This project presents RivetDB, a high-performance, concurrent in-memory database implemented in Rust that eliminates this trade-off. RivetDB supports over 70 commands across 10 data types, introduces a built-in SQL-like query engine, unlimited namespace-based multi-tenancy, and asynchronous AOF persistence. Built on the Tokio async runtime with DashMap for concurrent access, RivetDB achieves 43\% faster atomic operations than Redis while providing compile-time guaranteed memory safety.

% ============================================================
% LIST OF FIGURES
% ============================================================
\newpage
\vspace*{-2\baselineskip}
\listoffigures

% ============================================================
% LIST OF TABLES
% ============================================================
\newpage
\vspace*{-2\baselineskip}
\listoftables

% ============================================================
% TABLE OF CONTENTS
% ============================================================
\newpage
\vspace*{-2\baselineskip}
\tableofcontents

% ============================================================
% CHAPTER 1: INTRODUCTION
% ============================================================
\chapter{Introduction}

\section{Background}

In-memory databases (IMDBs) have become indispensable in modern computing environments that demand real-time data access, such as financial trading systems, e-commerce platforms, gaming servers, and IoT applications [1]. By storing all operational data in the main memory, IMDBs eliminate the latency overhead associated with disk I/O, enabling response times measured in microseconds. However, as the disparity between CPU processing speed and memory access latency continues to grow, modern database architectures must be rethought around memory-centric computing paradigms rather than traditional I/O optimization [2].

The two dominant in-memory systems---Redis and Memcached---represent opposite poles of an architectural trade-off rooted in their shared implementation language, C. Redis processes commands in a primarily single-threaded event loop, providing data consistency and a rich feature set at the cost of multi-core utilization. Memcached employs multi-threading for superior throughput but supports only simple key-value pairs with no persistence, severely limiting its applicability. Both systems remain susceptible to memory safety vulnerabilities inherent to C, including buffer overflows, use-after-free errors, and data races [9].

Furthermore, several capabilities essential for modern applications---native time-series data handling, probabilistic data structures such as Bloom filters, SQL-like querying, and multi-tenant data isolation---are either absent from these core systems or available only through separately maintained, often commercially licensed modules.

The Rust programming language offers a new paradigm for systems programming. Its ownership model, enforced at compile time, guarantees memory safety and data-race freedom without garbage collection overhead [9]. This project presents RivetDB, a database that leverages Rust's safety guarantees and the Tokio asynchronous runtime to provide a concurrent, feature-rich IMDB that unifies the strengths of both Redis and Memcached while introducing capabilities neither system offers.

\section{Relevance}

Modern applications across e-commerce, financial services, and the Internet of Things (IoT) increasingly depend on real-time data processing. This has made in-memory databases a critical infrastructure component [1]. However, the most widely adopted IMDBs present developers with a fundamental architectural trade-off rooted in the limitations of C [9]. There is no mainstream IMDB that successfully unifies a rich feature set, multi-core scalability, built-in advanced data types, and guaranteed safety in a single system. RivetDB addresses this gap.

\section{Project Undertaken}

The objective of this project is to design and implement \textbf{RivetDB}, a high-performance, concurrent, in-memory database system written in Rust that:

\begin{enumerate}[leftmargin=2.5em]
    \item Implements a Redis-compatible protocol (RESP) enabling seamless integration with existing Redis client libraries.
    \item Provides concurrent architecture using DashMap for true multi-threaded command execution.
    \item Supports 10 native data types including strings, lists, sets, sorted sets, hashes, JSON, Bloom filters, and time-series.
    \item Includes a built-in SQL-like query engine absent from standard Redis.
    \item Implements unlimited namespace-based multi-tenancy with per-tenant memory isolation.
    \item Provides asynchronous AOF persistence with configurable fsync policies.
    \item Incorporates authentication and security measures.
    \item Demonstrates capabilities through real-world application demos.
\end{enumerate}

\section{Organization of Project Report}

The remainder of this report is organized as follows:

\begin{itemize}[leftmargin=2.5em]
    \item \textbf{Chapter 2: Literature Survey} surveys existing research and systems related to in-memory databases, concurrency control, and memory-safe systems programming.
    \item \textbf{Chapter 3: System Design} describes the architecture, data flow, and design decisions of RivetDB.
    \item \textbf{Chapter 4: Implementation} details the technical implementation of core modules and features.
    \item \textbf{Chapter 5: Performance Optimization} describes the benchmark-driven optimization methodology.
    \item \textbf{Chapter 6: Results and Discussion} presents benchmark results, performance analysis, and comparisons with Redis.
    \item \textbf{Chapter 7: Conclusion and Future Work} summarizes findings and outlines directions for future enhancement.
\end{itemize}

% ============================================================
% CHAPTER 2: LITERATURE SURVEY
% ============================================================
\chapter{Literature Survey}

\section{Introduction}

The landscape of in-memory data storage is dominated by systems that have historically prioritized raw performance, often at the cost of safety or feature richness. This chapter surveys the foundational C-based systems that motivate our work, explores relevant academic research in concurrency, and examines the emerging ecosystem of databases built with modern systems languages like Rust.

\section{In-Memory Database Architecture}

A comprehensive overview of main-memory database systems established that eliminating disk I/O fundamentally changes design considerations for index structures, buffer management, and logging strategies [1]. More recent research has identified the ``memory wall''---the growing disparity between CPU speed and memory latency---and argues that databases must optimize for memory access patterns, not merely avoid disk [2]. Research into replicated data layouts for in-memory systems has demonstrated that organizing data to minimize cache coherence overhead across cores yields significant performance improvements [3]. These findings motivated RivetDB's use of DashMap's sharded architecture, where each shard operates on independent data to reduce cross-core contention.

\section{The C-Based Incumbents: Redis and Memcached}

For over a decade, the primary choice for in-memory data stores has been a trade-off between Redis and Memcached.

\textbf{Redis} is a feature-rich database known for its support of advanced data structures like lists, sets, and hashes, alongside persistence mechanisms such as append-only file (AOF) logging. However, to avoid concurrency-related bugs, its core command processing is primarily single-threaded. This design choice simplifies development and ensures data consistency but creates a significant performance bottleneck on modern multi-core processors.

\textbf{Memcached}, in contrast, was designed for maximum throughput and scalability through a multi-threaded architecture. It excels in distributed caching scenarios where raw speed is paramount. However, this performance comes with significant limitations: it supports only simple key-value pairs and offers no persistence mechanisms.

Both systems, being implemented in C, are susceptible to memory safety vulnerabilities such as buffer overflows and use-after-free errors, which can compromise system reliability and security [9].

\section{Concurrency Control}

A foundational survey of concurrency control categorizes approaches into pessimistic (locking-based) and optimistic (validation-based) families, analyzing trade-offs between lock overhead and abort rates [4]. Further investigation has demonstrated that neither approach is universally superior: optimistic methods perform well under low contention, while pessimistic approaches are preferable when conflicts are frequent [5]. Research on lock-free and wait-free data structures has shown that concurrent threads can make progress without costly locking mechanisms [6]. RivetDB employs a hybrid strategy: DashMap uses per-shard read-write locks providing optimistic-like performance under low contention, while atomic counters ensure observability paths never block command execution.

\section{Asynchronous I/O and Networking}

Research into NIC-software co-design has demonstrated that minimizing data copies and context switches in the I/O path yields significant throughput improvements [7]. Studies on asynchronous frameworks have shown that event-driven, non-blocking architectures handle orders of magnitude more concurrent connections than thread-per-connection models [8]. RivetDB builds on these principles through the Tokio async runtime, where each client connection is a lightweight task rather than a system thread.

\section{The Rust Paradigm: Safety and Fearless Concurrency}

A formal treatment of safe systems programming in Rust establishes that Rust's type system and ownership model provide provable guarantees about memory safety and data-race freedom without runtime overhead [9]. The key features relevant to database development include:

\begin{itemize}[leftmargin=2.5em]
    \item \textbf{Memory Safety}: Rust's ownership model and borrow checker guarantee memory safety without needing a garbage collector, statically preventing segmentation faults and memory corruption [9].
    \item \textbf{Fearless Concurrency}: The language enforces thread safety at compile time through its Send and Sync traits, preventing data races [9].
    \item \textbf{Performance}: Rust's zero-cost abstractions ensure that safety guarantees impose no runtime performance penalty.
\end{itemize}

Practical experience building Redis-like databases in Rust has established the viability of the platform and identified key challenges in memory management, connection handling, and persistence [10]. RivetDB extends beyond these educational foundations with concurrent data structures, a SQL query engine, multi-tenancy, and 10 native data types.

% ============================================================
% CHAPTER 3: SYSTEM DESIGN
% ============================================================
\chapter{System Design}

\section{Introduction}

This chapter describes the architectural design of RivetDB, detailing the modular system structure and the rationale behind key design decisions.

\section{Architectural Overview}

RivetDB is designed as a modular system with five primary layers:

\begin{enumerate}[leftmargin=2.5em]
    \item \textbf{Networking Layer}: Tokio async runtime with TCP listener and per-connection authentication state tracking.
    \item \textbf{Command Router}: 70+ commands dispatched across 16 modules.
    \item \textbf{Storage Engine}: DashMap-based concurrent hash map with 10 native data types.
    \item \textbf{Advanced Features Layer}: SQL query engine, namespace-based multi-tenancy, and Prometheus metrics.
    \item \textbf{Persistence Layer}: Asynchronous AOF via crossbeam channels with configurable fsync policies.
\end{enumerate}

\begin{figure}[htbp]
\centering
\begin{verbatim}
   Client Applications (Any Redis Client)
              |
              | RESP Protocol (Redis-compatible)
              v
   +--------------------------------------+
   |       Networking Layer               |
   |  Tokio Async Runtime + TCP           |
   |  Authentication (AUTH command)       |
   +--------------------------------------+
   |       Command Router                 |
   |  70+ commands across 16 modules      |
   +--------------------------------------+
   |       Storage Engine                 |
   |  DashMap (sharded concurrent map)    |
   |  10 native data types               |
   +--------------------------------------+
   |     Advanced Features Layer          |
   |  SQL Query | Multi-Tenancy | Metrics |
   +--------------------------------------+
   |       Persistence Layer              |
   |  Async AOF via crossbeam channel     |
   |  Configurable fsync policies         |
   +--------------------------------------+
\end{verbatim}
\caption{RivetDB System Architecture}
\label{fig:architecture}
\end{figure}

\section{Networking and Protocol Layer}

RivetDB implements the Redis Serialization Protocol (RESP), enabling compatibility with all existing Redis client libraries across every major programming language. The networking layer is built on Tokio's async runtime [8], where each client connection is handled as a lightweight asynchronous task. This design enables thousands of concurrent connections with minimal memory overhead compared to thread-per-connection models.

Connection handling includes per-connection authentication state tracking. When security is enabled, clients must issue an AUTH command before any other commands are accepted, with exempt commands (PING, QUIT) allowed for health checking.

\section{Storage Engine}

The core storage engine uses DashMap---a concurrent hash map with per-shard read-write locks---pre-sized to 100,000 entries to minimize rehashing under load. DashMap divides data across multiple shards, each independently lockable, enabling true parallel read and write access from multiple threads [5][6].

RivetDB natively supports 10 data types as shown in Table~\ref{tab:datatypes}.

\begin{table}[htbp]
\caption{RivetDB Native Data Types and Commands}
\label{tab:datatypes}
\centering
\small
\begin{tabular}{@{}p{2.5cm}p{4cm}p{5cm}@{}}
\toprule
\textbf{Data Type} & \textbf{Key Commands} & \textbf{Description} \\
\midrule
Strings & SET, GET, INCR, DECR, MGET, MSET, APPEND, STRLEN & Basic key-value with numeric operations \\
\midrule
Lists & LPUSH, RPUSH, LPOP, RPOP, LLEN, LRANGE, LINDEX, LSET & Ordered collections \\
\midrule
Sets & SADD, SREM, SMEMBERS, SISMEMBER, SUNION, SINTER, SDIFF & Unique element collections \\
\midrule
Sorted Sets & ZADD, ZRANGE, ZREVRANGE, ZRANGEBYSCORE, ZSCORE, ZINCRBY & Score-ordered rankings \\
\midrule
Hashes & HSET, HGET, HMGET, HGETALL, HDEL, HINCRBY, HKEYS, HVALS & Field-value maps \\
\midrule
JSON & JSON.SET, JSON.GET, JSON.DEL, JSON.ARRAPPEND, JSON.OBJKEYS & Document storage \\
\midrule
Bloom Filters & BF.RESERVE, BF.ADD, BF.EXISTS, BF.INFO, BF.CARD & Probabilistic membership \\
\midrule
Time-Series & TS.CREATE, TS.ADD, TS.RANGE, TS.MRANGE, TS.INFO, TS.INCRBY & Timestamped data \\
\midrule
Schemas & TSET, TGET, TKEYS, TVALIDATE & Typed storage \\
\midrule
Namespaces & NAMESPACE CREATE/ USE/ DELETE/ LIST/ STATS & Tenant isolation \\
\bottomrule
\end{tabular}
\end{table}

\section{SQL Query Engine}

RivetDB includes a built-in SQL-like query engine---a capability absent from standard Redis entirely. The engine supports:

\begin{lstlisting}
QUERY "SELECT * FROM keys WHERE type = 'string'"
QUERY "SELECT * FROM keys WHERE key LIKE 'user:*'"
QUERY "SELECT COUNT(*) FROM keys WHERE type = 'timeseries'"
QUERY "DELETE FROM keys WHERE key LIKE 'temp:*'"
EXPLAIN "SELECT * FROM keys WHERE type = 'hash'"
\end{lstlisting}

The query executor parses SQL statements, translates them into key-space scans with predicate filtering, and returns tabular results. The EXPLAIN command provides query plan visualization for optimization.

\section{Multi-Tenancy via Namespaces}

Unlike Redis, which limits users to 16 numbered databases without memory isolation, RivetDB implements namespace-based multi-tenancy with unlimited namespaces, per-tenant memory limits, usage statistics tracking, and complete data isolation between tenants.

\section{Persistence Layer}

RivetDB implements AOF persistence with three fsync policies: \textbf{Always} (flush after every write), \textbf{Everysec} (buffer up to 1,000 writes), and \textbf{No} (buffer up to 10,000 writes). The asynchronous AOF writer uses crossbeam channels [8] to decouple disk I/O from command processing, with a 64KB BufWriter (8$\times$ default) to minimize system call overhead.

\section{Observability and Security}

RivetDB provides a Prometheus-compatible metrics endpoint, slow query logging (commands exceeding 1ms), per-command statistics via DashMap counters, and hot key detection. Authentication is implemented via a Redis-compatible AUTH command with configurable password and per-connection state tracking.

% ============================================================
% CHAPTER 4: IMPLEMENTATION
% ============================================================
\chapter{Implementation}

\section{Introduction}

This chapter details the phased implementation of RivetDB, from the initial prototype through advanced feature integration.

\section{Phase 1: Core Key-Value Store Prototype}

Development began with a minimal viable system implementing basic key-value operations (GET, SET, DELETE). A TCP server was built using Tokio to handle client requests [10], validating the fundamental design and its ability to manage concurrent connections.

\section{Phase 2: Asynchronous and Concurrent Architecture}

The system was built on Rust's asynchronous runtime, Tokio, to handle network I/O without blocking threads. Thread-safe primitives, primarily \texttt{Arc<DashMap<K,V>{}>}, were used to wrap the core data store [5]. Data is sharded across multiple independent stores, with each shard managed independently, allowing requests to be processed in parallel with minimal lock contention. Rust's compile-time safety guarantees were crucial in this stage, preventing data races by design [9].

\section{Phase 3: Feature Expansion}

With a stable concurrent core, the system was extended to provide a feature-rich data model. This included implementing:

\begin{itemize}[leftmargin=2.5em]
    \item \textbf{Lists, Sets, Sorted Sets, Hashes}: Core Redis-compatible data structures with full command coverage.
    \item \textbf{JSON Documents}: Structured document storage with path-based access, array operations, and numeric increments.
    \item \textbf{Bloom Filters}: Built-in probabilistic membership testing with configurable error rates.
    \item \textbf{Time-Series}: Native timestamped data ingestion with aggregation, retention policies, and label-based querying.
    \item \textbf{SQL Query Engine}: SELECT, DELETE, COUNT, and EXPLAIN operations over the key-space.
    \item \textbf{Namespaces}: Unlimited tenant-isolated data spaces with per-tenant memory limits.
\end{itemize}

\section{Phase 4: Persistence}

An append-only file (AOF) mechanism was implemented, logging every write operation to disk for crash recovery [10]. The key innovation is the asynchronous AOF writer using crossbeam channels, which completely decouples persistence from command execution.

\section{Phase 5: Authentication and Observability}

Authentication via the AUTH command was added with per-connection state tracking. A Prometheus-compatible metrics endpoint was implemented, alongside slow query logging and command statistics.

% ============================================================
% CHAPTER 5: PERFORMANCE OPTIMIZATION
% ============================================================
\chapter{Performance Optimization}

\section{Introduction}

Performance optimization followed a benchmark-driven, phased approach using \texttt{redis-benchmark} as the industry-standard measurement tool.

\section{Phase 1: Quick Wins}

\begin{itemize}[leftmargin=2.5em]
    \item \textbf{BufWriter capacity}: Increased from 8KB (default) to 64KB, reducing system call frequency for sequential writes.
    \item \textbf{DashMap pre-sizing}: Initialized with capacity for 100,000 entries, eliminating rehash pauses during high-volume ingestion.
\end{itemize}

\textbf{Result}: $\sim$22\% improvement in pipelined SET throughput, $\sim$18\% latency reduction.

\section{Phase 2: Asynchronous AOF}

\begin{itemize}[leftmargin=2.5em]
    \item Implemented crossbeam channel-based async writer, decoupling disk I/O from command processing.
    \item Increased Everysec flush threshold from 100 to 1,000 writes.
    \item Increased No-policy flush threshold from 1,000 to 10,000 writes.
\end{itemize}

\textbf{Result}: Additional $\sim$78\% improvement in pipelined SET throughput (cumulative $\sim$100\% from baseline).

% ============================================================
% CHAPTER 6: RESULTS AND DISCUSSION
% ============================================================
\chapter{Results and Discussion}

\section{Introduction}

This chapter presents the experimental results of benchmarking RivetDB against Redis, followed by analysis and discussion of the findings.

\section{Experimental Setup}

\begin{itemize}[leftmargin=2.5em]
    \item \textbf{Hardware}: Consumer-grade laptop (representative of development environments).
    \item \textbf{Software}: RivetDB compiled with \texttt{--release} optimizations; Redis latest stable via Docker.
    \item \textbf{Tool}: \texttt{redis-benchmark} (industry standard), ensuring identical test conditions.
    \item \textbf{Configuration}: AOF enabled with Everysec policy on both systems.
    \item \textbf{Parameters}: 10,000 requests per test, 50 concurrent connections.
\end{itemize}

\section{Throughput Results}

\subsection{Standard (Non-Pipelined) Operations}

\begin{table}[htbp]
\caption{Standard (Non-Pipelined) Throughput Comparison}
\label{tab:standard}
\centering
\begin{tabular}{@{}lrrrr@{}}
\toprule
\textbf{Command} & \textbf{RivetDB} & \textbf{Redis} & \textbf{Ratio} & \textbf{Winner} \\
\midrule
INCR & 14,741 req/s & 10,322 req/s & 143\% & RivetDB \\
GET & $\sim$14,000 req/s & 12,262 req/s & $\sim$114\% & RivetDB \\
\bottomrule
\end{tabular}
\end{table}

In standard mode, where each command completes before the next is sent, RivetDB outperforms Redis on both read and write operations (Table~\ref{tab:standard}). The INCR command---a read-modify-write atomic operation commonly used in counters, rate limiters, and metrics---is 43\% faster on RivetDB. Standard GET operations show a 14\% advantage, demonstrating that DashMap's concurrent read access is efficient even for single-command round trips.

\subsection{Pipelined Operations}

\begin{table}[htbp]
\caption{Pipelined Throughput Comparison (16$\times$ Pipeline Depth)}
\label{tab:pipelined}
\centering
\begin{tabular}{@{}lrrrr@{}}
\toprule
\textbf{Command} & \textbf{RivetDB} & \textbf{Redis} & \textbf{Ratio} & \textbf{Winner} \\
\midrule
GET & 104,384 req/s & 155,763 req/s & 67\% & Redis \\
SET & 8,846 req/s & 145,985 req/s & 6\% & Redis \\
\bottomrule
\end{tabular}
\end{table}

Under aggressive 16$\times$ pipelining (Table~\ref{tab:pipelined}), Redis's single-threaded event loop processes batches sequentially without lock overhead. The SET gap is primarily attributable to AOF persistence overhead.

\section{Latency Results}

\begin{table}[htbp]
\caption{Median Latency Comparison (p50)}
\label{tab:latency}
\centering
\begin{tabular}{@{}lrrl@{}}
\toprule
\textbf{Command} & \textbf{RivetDB} & \textbf{Redis} & \textbf{Improvement} \\
\midrule
INCR & 2.43 ms & 4.12 ms & 41\% lower \\
GET (pipelined) & 5.32 ms & 4.34 ms & Comparable \\
SET (pipelined) & 78.98 ms & 4.81 ms & AOF bottleneck \\
\bottomrule
\end{tabular}
\end{table}

RivetDB achieves 41\% lower median latency than Redis on INCR operations (Table~\ref{tab:latency}).

\section{Optimization Impact}

\begin{table}[htbp]
\caption{Performance Before and After Optimization}
\label{tab:optprogress}
\centering
\begin{tabular}{@{}lrrl@{}}
\toprule
\textbf{Metric} & \textbf{Baseline} & \textbf{After Phase 1+2} & \textbf{Change} \\
\midrule
SET (pipelined req/s) & 7,252 & 8,846 & +22\% \\
SET latency (p50) & 78.98 ms & 65.06 ms & $-$18\% \\
GET (pipelined req/s) & 104,384 & 95,511 & $\sim$same \\
INCR (req/s) & 14,741 & 14,575 & $\sim$same \\
\bottomrule
\end{tabular}
\end{table}

\begin{table}[htbp]
\caption{Post-Optimization Comparison vs Redis}
\label{tab:postopt}
\centering
\begin{tabular}{@{}lrrl@{}}
\toprule
\textbf{Command} & \textbf{RivetDB} & \textbf{Redis} & \textbf{Ratio} \\
\midrule
INCR & 14,575 req/s & 10,865 req/s & 134\% \\
GET (pipelined) & 95,511 req/s & 156,495 req/s & 61\% \\
SET (pipelined) & 8,846 req/s & 145,985 req/s & 6\% \\
\bottomrule
\end{tabular}
\end{table}

Even after optimization (Tables~\ref{tab:optprogress} and~\ref{tab:postopt}), RivetDB maintains its 34\% lead in atomic operations, confirming that DashMap's fine-grained concurrent access is fundamentally more efficient than Redis's single-threaded model for non-batched operations.

\section{Estimated Performance Without AOF}

\begin{table}[htbp]
\caption{Estimated Performance Without AOF Persistence}
\label{tab:noaof}
\centering
\begin{tabular}{@{}lrr@{}}
\toprule
\textbf{Command} & \textbf{With AOF} & \textbf{Without AOF (Est.)} \\
\midrule
SET (pipelined) & 8,846 req/s & $\sim$60,000--80,000 req/s \\
SET (standard) & $\sim$14,000 req/s & $\sim$14,500--15,000 req/s \\
SET latency (p50) & 65.06 ms & $\sim$3--5 ms \\
\bottomrule
\end{tabular}
\end{table}

AOF-related operations consume approximately 85--90\% of the SET execution path (Table~\ref{tab:noaof}). For applications where durability is not required (caching, session stores), disabling AOF would bring RivetDB's write throughput to within 40--55\% of Redis.

\section{Feature Comparison}

\begin{table}[htbp]
\caption{Feature Completeness: RivetDB vs Redis (Standard)}
\label{tab:features}
\centering
\begin{tabular}{@{}lll@{}}
\toprule
\textbf{Feature} & \textbf{Redis} & \textbf{RivetDB} \\
\midrule
Core data types & 5 & 10 \\
Time-Series & Paid module & Built-in \\
SQL Query & Not available & Built-in \\
Bloom Filters & Paid module & Built-in \\
JSON documents & Paid module & Built-in \\
Multi-Tenancy & 16 fixed DBs & Unlimited \\
Memory safety & Manual (C) & Compile-time \\
Prometheus metrics & Not built-in & Built-in \\
\bottomrule
\end{tabular}
\end{table}

\section{Discussion}

RivetDB's performance profile reveals a clear architectural trade-off. RivetDB excels in standard (non-pipelined) operations where its concurrent architecture processes independent requests in parallel across DashMap shards. The 43\% INCR and 14\% GET advantage demonstrates that multi-threaded access patterns outperform single-threaded event loops for independent operations. Redis excels in pipelined operations where batched commands benefit from sequential execution without lock acquisition overhead.

RivetDB's decisive advantage is feature completeness (Table~\ref{tab:features}). Five advanced capabilities that require commercially licensed Redis modules---or are entirely absent---are natively integrated into RivetDB's single binary.

\subsection{The Impact of Rust's Safety Model}

Rust's ownership model and borrow checker fundamentally changed the development experience compared to C-based database development [9]. The compiler's enforcement of thread safety at compile time eliminated data races, use-after-free errors, and null pointer dereferences during development. DashMap, Tokio, and crossbeam---all libraries that internally use Rust's unsafe mechanisms---provide safe abstractions verified through the RustBelt framework [9]. RivetDB's application code contains no \texttt{unsafe} blocks.

\subsection{Application Demonstrations}

Two real-world demonstrations validate RivetDB's capabilities:

\textbf{Collaborative Whiteboard}: A multi-user drawing application with real-time cursor synchronization across networked machines, exercising HSET (cursors), LPUSH (strokes), and SADD (online users) simultaneously.

\textbf{Gaming Leaderboard}: A real-time ranking system using ZADD/ZREVRANGE (sorted sets), TS.ADD/TS.RANGE (time-series), and JSON.SET/JSON.GET (player profiles)---three features requiring separate Redis modules.

% ============================================================
% CHAPTER 7: CONCLUSION AND FUTURE WORK
% ============================================================
\chapter{Conclusion and Future Work}

\section{Conclusion}

This project successfully demonstrated the design and implementation of RivetDB, a high-performance, concurrent in-memory database in Rust that overcomes the long-standing architectural trade-offs inherent in C-based systems like Redis and Memcached [10]. The key contributions are:

\begin{enumerate}[leftmargin=2.5em]
    \item \textbf{Multi-threaded feature-rich architecture}: RivetDB supports 70+ commands across 10 native data types with true concurrent access via DashMap, eliminating Redis's single-threaded limitation [10].

    \item \textbf{Novel built-in capabilities}: A SQL-like query engine, unlimited namespace-based multi-tenancy, and native time-series, Bloom filter, and JSON support---features unavailable in standard Redis---are integrated into a single binary.

    \item \textbf{Benchmark-validated performance}: RivetDB achieves 14,741 atomic ops/sec (43\% faster than Redis) and 41\% lower INCR latency while providing compile-time guaranteed safety [9].

    \item \textbf{Asynchronous persistence}: The crossbeam channel-based AOF writer decouples disk I/O from command execution. Phased optimizations achieved a cumulative 100\% improvement in write throughput.

    \item \textbf{Practical validation}: Real-world application demos---a collaborative whiteboard and gaming leaderboard---demonstrate suitability for latency-sensitive applications.
\end{enumerate}

RivetDB establishes that modern systems languages can redefine the design space for high-performance databases, proving that multi-threaded performance, rich features, and memory safety are not mutually exclusive goals [9].

\section{Future Work}

\begin{itemize}[leftmargin=2.5em]
    \item \textbf{Distributed Clustering}: Implement horizontal scaling through data sharding across nodes, with consensus algorithms (e.g., Raft) for fault tolerance [4].
    \item \textbf{Write Path Optimization}: Implement batch command processing to accumulate pipelined writes into a single AOF entry, targeting Redis-competitive pipelined SET performance.
    \item \textbf{TLS Encryption}: Integrate rustls for encrypted connections to support production deployment.
    \item \textbf{Access Control Lists}: Implement per-user command and key permissions beyond basic password authentication.
    \item \textbf{Pub/Sub Messaging}: Add publish/subscribe channels for event-driven architectures.
    \item \textbf{RDB Snapshots}: Implement point-in-time snapshot persistence as a complement to AOF.
\end{itemize}

\subsection{Industrial Relevance}

The industrial relevance of a safe, scalable, and feature-rich IMDB is significant. Such a system is ideally suited for mission-critical applications in financial services, e-commerce platforms, real-time gaming servers, and IoT data ingestion pipelines [1].

% ============================================================
% REFERENCES
% ============================================================
\begin{thebibliography}{10}

\bibitem{ref1}
P. Larson and J. Levandoski, ``An Overview of Main-Memory Database Systems,'' 2006. [Online]. Available: \url{https://www.vldb.org/pvldb/vol9/p1609-larson.pdf}

\bibitem{ref2}
Y. Chronis \textit{et al.}, ``The Memory Wall and The Case for Memory-Centric Database Systems,'' in \textit{CIDR 2025}, 2025. [Online]. Available: \url{https://vldb.org/cidrdb/papers/2025/p6-chronis.pdf}

\bibitem{ref3}
S. Sudhir, M. Cafarella, and S. Madden, ``Replicated Layout for In-Memory Database Systems,'' \textit{PVLDB}, vol. 15, no. 4, pp. 984--997, 2022. [Online]. Available: \url{https://www.vldb.org/pvldb/vol15/p984-sudhir.pdf}

\bibitem{ref4}
P. A. Bernstein and N. Goodman, ``Concurrency Control in Distributed Database Systems,'' \textit{ACM Comput. Surv.}, 1981. [Online]. Available: \url{https://people.eecs.berkeley.edu/~brewer/cs262/concurrency-distributed-databases.pdf}

\bibitem{ref5}
G. Graefe, ``Revisiting Optimistic and Pessimistic Concurrency Control,'' Hewlett Packard Labs, 2016. [Online]. Available: \url{https://www.labs.hpe.com/techreports/2016/HPE-2016-47.pdf}

\bibitem{ref6}
E. Petrank and S. Timnat, ``A Practical Wait-Free Simulation for Lock-Free Data Structures,'' in \textit{PPoPP 2014}, 2014. [Online]. Available: \url{https://csaws.cs.technion.ac.il/~erez/Papers/wf-simulation-ppopp14.pdf}

\bibitem{ref7}
A. Skiadopoulos \textit{et al.}, ``ZeroNIC: A NIC-Software Co-Design for High-Performance and Flexible Host Networking,'' in \textit{OSDI 2024}, 2024. [Online]. Available: \url{https://www.usenix.org/system/files/osdi24-skiadopoulos.pdf}

\bibitem{ref8}
As if, D. Yadav, and A. Yadav, ``Using Asynchronous Frameworks and Database Connection Pools to Enhance Web Application Performance in High-Concurrency Environments,'' 2024. [Online]. Available: \url{https://www.researchgate.net/publication/385174027}

\bibitem{ref9}
R. Jung, J. Jourdan, R. Krebbers, and D. Dreyer, ``Safe Systems Programming in Rust,'' \textit{Communications of the ACM}, 2021. [Online]. Available: \url{https://cacm.acm.org/research/safe-systems-programming-in-rust/}

\bibitem{ref10}
A. Kumar, ``Building a Redis-like In-Memory Database from Scratch in Rust,'' \textit{Medium}, 2023. [Online]. Available: \url{https://medium.com/rustaceans/my-journey-into-rust-building-a-redis-like-in-memory-database-from-scratch-a622c755065d}

\end{thebibliography}

\end{document}
